\documentclass[../DoAn.tex]{subfiles}
\begin{document}

% =====================================================================
% CHƯƠNG 7 — KẾT QUẢ THỰC NGHIỆM VÀ ĐÁNH GIÁ
% "Results & Discussion" — nơi ăn điểm "hàm lượng khoa học".
% =====================================================================

\section{Giới thiệu chương}

Chương này trình bày và phân tích kết quả thực nghiệm của mô hình đề xuất KAN-WNETR nhằm trả lời câu hỏi cốt lõi: việc thay khối MLP/FFN cố định bằng KAN có thực sự cải thiện chất lượng trích xuất fECG hay không, và mô hình có đủ nhẹ để triển khai trên thiết bị biên hay không. Để đánh giá toàn diện, kết quả được tổ chức theo cả hai khía cạnh định tính và định lượng. Trước hết, phần kết quả định tính khảo sát trực quan dạng sóng fECG tái tạo ở nhiều mức nhiễu. Tiếp theo, chất lượng tín hiệu được lượng hóa bằng SSIM/PSNR và đối sánh với W-NETR, sau đó là hiệu quả phát hiện phức bộ fQRS và so sánh với các phương pháp tiên tiến (SOTA). Cuối cùng, chương phân tích độ phức tạp tính toán làm cầu nối sang phần triển khai hệ thống biên (Chương~\ref{chapter:System}), rồi bàn luận về điểm mạnh, hạn chế và đánh đổi giữa hai phiên bản FP32 và INT8.

% ---------------------------------------------------------------------
\section{Kết quả định tính}

Hình~\ref{fig:result} biểu diễn kết quả của mô hình đề xuất trên bộ dữ liệu FECGSYNDB. Tín hiệu được hiển thị là ở C3 của tín hiệu maECG trên 3 mức nhiễu SNR 0dB, 6dB và 12dB. Ta có thể thấy rằng bất cứ trường hợp nhiễu nào, KAN-WNETR vẫn có thể phân tách được tín hiệu fECG, cho thấy hiệu năng cao của mô hình này. Ngoài bắt được các phức hợp đỉnh QRS hoàn chỉnh, mô hình đồng thời cũng hiển thị lại được rõ các sóng P và T, là những sóng được đánh giá là khó để thu lại được \cite{leeEndtoEndDeepLearning2022}.

\begin{figure}[h]
  \centering
  \includegraphics[width=1\textwidth]{Hinhve/image_46.png}
  \caption{Kết quả trích xuất fECG của KAN-WNETR trên FECGSYNDB (case C3, SNR 0/6/12 dB).}
  \label{fig:result}
\end{figure}

% ---------------------------------------------------------------------
\section{Chất lượng tín hiệu định lượng (SSIM/PSNR)}

Để phân tích cụ thể hơn, kết quả tham số PSNR và SSIM đánh giá chất lượng tín hiệu dựng lại được thể hiện trong bảng \ref{tab:ssim_compare_wnetr}, đồng thời mô hình được so sánh với kiến trúc W-NETR \cite{almadaniOneDimensionalWNETRNonInvasive2023}, được cho là có kết quả tốt nhất vào thời điểm viết bài. W-NETR đạt SSIM tổng là $0.96$ (PSNR tổng $51.8$ dB) khi đánh giá trên tất cả các loại nhiễu C0–C5.

\begin{table}[htbp]
\centering
\caption{So sánh SSIM theo từng loại nhiễu trên FECGSYNDB (C0--C5).}
\label{tab:ssim_compare_wnetr}
\begin{tabular}{lccccccc}
\hline
\textbf{Mô hình} & \textbf{C0} & \textbf{C1} & \textbf{C2} & \textbf{C3} & \textbf{C4} & \textbf{C5} & \textbf{Total} \\
\hline
W-NETR (paper)      & \textbf{0.990} & \textbf{0.990} & \textbf{0.990} & 0.910 & 0.930 & 0.970 & 0.960 \\
Our KAN-WNETR & 0.976 & 0.972 & 0.973 & \textbf{0.933} & \textbf{0.976} & \textbf{0.975} & \textbf{0.966} \\
\hline
\end{tabular}
\end{table}

Kết quả SSIM của mô hình KAN-WNETR trên tập test cho từng case (c0–c5) được tổng hợp và đối sánh với W-NETR như Bảng~\ref{tab:ssim_compare_wnetr}. Nhìn chung, SSIM trung bình của mô hình đề xuất đạt xấp xỉ $0.963$, tương đương mức SSIM tổng $0.96$ mà W-NETR báo cáo. Tuy nhiên, phân rã theo từng loại nhiễu cho thấy xu hướng khác biệt: mô hình của chúng tôi thể hiện lợi thế rõ hơn ở các trường hợp nhiễu “khó” (đặc biệt C3–C4), trong khi các case (C0–C2) vẫn thấp hơn so với W-NETR.

Các kết quả này gợi ý rằng cấu hình KAN-WNETR có thể mang lại lợi thế về “bảo toàn cấu trúc” trong điều kiện nhiễu/biến thiên khó, nhưng vẫn cần tối ưu thêm để tiệm cận trần hiệu năng ở các kịch bản đơn giản.

% ---------------------------------------------------------------------
\section{Phát hiện fQRS và so sánh SOTA}

\begin{table}[htbp]
\centering
\caption{So sánh hiệu quả với các phương pháp SOTA (Recall, Precision, F1).}
\label{tab:compare_recall_precision_f1}
\setlength{\tabcolsep}{8pt}
\renewcommand{\arraystretch}{1.15}
\begin{tabular}{lccc}
\hline
\textbf{Method} & \textbf{Recall (\%)} & \textbf{Precision (\%)} & \textbf{F1 (\%)} \\
\hline
(EKF)                     & 85.58 & 89.53 & 86.96 \\
(TS$_{\mathrm{PCA}}$)     & 90.91 & 93.55 & 92.17 \\
RCED-Net                  & 92.69 & 95.67 & 93.73 \\
Res-Unet                  & 91.88 & 94.36 & 93.10 \\
AECG-DecomposeNet         & 93.52 & 97.04 & 95.42 \\
Attention-based CycleGAN  & 96.70 & 95.90 & 96.30 \\
STFT GANs                 & --    & --    & 96.83 \\
W-net                     & 95.74 & 97.45 & 96.91 \\
W-NETR           & 98.12 & 98.60 & 98.90 \\
\hline
\textbf{Proposed KAN-WNETR (Ours)}   & \textbf{97.53} & \textbf{97.94} & \textbf{97.73} \\
\hline
\end{tabular}
\end{table}

Kết quả của mô hình KAN-WNETR đạt Recall = 97.53\%, Precision = 97.94\% và F1 = 97.73\% trên tập đánh giá. Với mức recall cao, mô hình có khả năng phát hiện phần lớn các đỉnh R của fECG mục tiêu (giảm bỏ sót), trong khi precision cao cho thấy số lượng phát hiện sai (false positives) được kiểm soát tốt. Việc F1 đạt xấp xỉ 97.7\% phản ánh sự cân bằng ổn định giữa hai tiêu chí này, đặc biệt quan trọng trong bối cảnh fECG có SNR thấp và dễ bị nhiễu/mECG che lấp. 

Đối chiếu KAN-WNETR với các phương pháp SOTA khác như EKF, $TS_{PCA}$, RCED-Net, Res-UNet, CycleGAN, W-net và W-NETR (Bảng~\ref{tab:compare_recall_precision_f1}), kết quả của mô hình đề xuất cao hơn rõ rệt so với nhóm phương pháp cổ điển, cho thấy lợi thế của hướng tiếp cận học sâu end-to-end trong bài toán tách và phát hiện fQRS. Mặc dù vẫn còn khoảng cách nhất định về F1 score khi so với kiến trúc W-NETR gốc, nhưng bù lại, KAN-WNETR lại nhỉnh hơn ở khía cạnh bảo toàn hình thái tín hiệu (SSIM) trong các ca nhiễu khó như đã phân tích ở phần trước.

% ---------------------------------------------------------------------
\section{Phân tích độ phức tạp}

Bên cạnh độ chính xác, khả năng triển khai thực tế là một tiêu chí quan trọng của đề tài, vốn hướng tới một hệ thống theo dõi đầu giường chạy trên thiết bị biên. Mô hình KAN-WNETR ở biến thể \texttt{kan\_wnetr\_base} có khoảng $38{,}87$ triệu tham số ở định dạng FP32. Do việc thay MLP/FFN bằng KAN chỉ tác động vào khối biến đổi phi tuyến theo token mà không làm thay đổi cơ chế self-attention hay macro-architecture chữ W, tổng chi phí tính toán của mô hình vẫn cùng bậc độ lớn với W-NETR gốc; phần phụ trội đến từ việc đánh giá các hàm cơ sở của KAN được kiểm soát nhờ sử dụng biến thể FasterKAN với hàm cơ sở RBF thay cho B-spline truyền thống.

Để phù hợp với ràng buộc tài nguyên của thiết bị biên, mô hình được chuyển sang định dạng ONNX và lượng tử hóa động INT8 cho các phép MatMul/Gemm. Sau lượng tử hóa, kích thước mô hình giảm còn khoảng $29$~MB, đủ nhẹ để nạp và chạy bằng ONNX Runtime trên CPU ARM của Raspberry Pi~5. Về tốc độ, mỗi lần suy luận trên một cửa sổ $992$ mẫu (tương đương $3{,}97$~giây tín hiệu) chỉ mất khoảng $80$--$120$~ms, tức nhanh hơn thời gian thực khoảng $8$--$12$ lần. Biên độ dự trữ tốc độ này cho phép hệ thống vừa suy luận liên tục theo cơ chế cửa sổ trượt, vừa còn dư tài nguyên cho thu nhận tín hiệu, hậu xử lý và truyền dữ liệu. Chi tiết cấu hình và số liệu benchmark được trình bày trong Bảng~\ref{tab:benchmark_int8} ở Chương~\ref{chapter:System}; phần đánh đổi chất lượng giữa FP32 và INT8 được bàn ở Mục~\ref{sec:discussion} ngay sau đây.

% ---------------------------------------------------------------------
\section{Bàn luận}
\label{sec:discussion}

Tổng hợp các kết quả trên, mô hình KAN-WNETR thể hiện một số \textbf{điểm mạnh} rõ rệt. Thứ nhất, ở khía cạnh bảo toàn hình thái, mô hình đạt SSIM tổng $0{,}966$ --- nhỉnh hơn mức $0{,}960$ mà W-NETR báo cáo --- và đặc biệt cải thiện ở các kịch bản nhiễu khó (C3--C4), nơi tín hiệu fECG dễ bị biến dạng hoặc che lấp nhất. Điều này phù hợp với giả thuyết ban đầu: hàm kích hoạt học được của KAN giúp khối biến đổi phi tuyến theo token thích nghi tốt hơn với các dạng méo phi tuyến trong tín hiệu sinh học. Thứ hai, về phát hiện fQRS, mô hình đạt sự cân bằng ổn định giữa Recall ($97{,}53\%$) và Precision ($97{,}94\%$), cho thấy khả năng vừa giảm bỏ sót vừa kiểm soát báo sai --- yếu tố then chốt khi ước lượng FHR đáng tin cậy trong môi trường SNR thấp.

Bên cạnh đó, đề tài cũng nhìn nhận thẳng thắn các \textbf{hạn chế}. Về độ chính xác phát hiện đỉnh, F1 của KAN-WNETR ($97{,}73\%$) vẫn còn một khoảng cách so với trần hiệu năng mà W-NETR gốc công bố ở một số thiết lập; lợi thế của mô hình đề xuất nghiêng về bảo toàn hình thái (SSIM) hơn là về chỉ số phát hiện thuần túy. Quan trọng hơn, phần lớn kết quả định lượng được kiểm chứng trên dữ liệu mô phỏng FECGSYNDB và một lượng hạn chế dữ liệu thực; do đó các con số này nên được hiểu là bằng chứng về tiềm năng phương pháp, chưa phải kết luận lâm sàng.

Cuối cùng, về \textbf{đánh đổi giữa FP32 và INT8}, việc lượng tử hóa động phục vụ triển khai biên tất yếu làm chất lượng tín hiệu suy giảm một mức nhất định so với mô hình FP32 gốc, do sai số biểu diễn trọng số ở độ phân giải 8-bit. Tuy nhiên, mức suy giảm này nhỏ và được xem là đánh đổi chấp nhận được khi cân nhắc lợi ích về kích thước mô hình (còn $\sim$$29$~MB) và tốc độ suy luận thời gian thực trên Raspberry Pi~5. Nói cách khác, bản INT8 hy sinh một phần nhỏ độ chính xác để đổi lấy tính khả thi triển khai --- yếu tố quyết định để mô hình rời khỏi phòng thí nghiệm và vận hành trong một hệ thống theo dõi thực tế.

% ---------------------------------------------------------------------
\section{Kết luận chương}

Chương này đã trình bày đầy đủ kết quả thực nghiệm của mô hình KAN-WNETR trên bài toán trích xuất fECG. Về chất lượng tái tạo dạng sóng, mô hình đạt SSIM tổng $0{,}966$, vượt nhẹ W-NETR và cải thiện rõ ở các trường hợp nhiễu khó. Về phát hiện phức bộ fQRS, mô hình đạt Recall $97{,}53\%$, Precision $97{,}94\%$ và F1 $97{,}73\%$, thuộc nhóm dẫn đầu khi đối sánh với các phương pháp tiên tiến. Đồng thời, phân tích độ phức tạp cho thấy mô hình sau lượng tử hóa INT8 đủ nhẹ và đủ nhanh để chạy thời gian thực trên thiết bị biên. Những kết quả này khẳng định tính khả thi của hướng tiếp cận thay MLP/FFN bằng KAN, đồng thời tạo tiền đề cho việc hiện thực hóa mô hình thành một hệ thống theo dõi hoàn chỉnh --- nội dung được trình bày trong Chương~\ref{chapter:System}.

\end{document}
